\chapter*{مقدمه: چرا این کتاب؟}
\addcontentsline{toc}{chapter}{مقدمه}

\begin{tcolorbox}[
    enhanced,
    colback=lightbg,
    colframe=darkgray,
    boxrule=1pt,
    arc=5pt
]
\begin{center}
\textit{«اگر نمی‌دانی از کجا آمده‌ای، نمی‌دانی کجا می‌روی.»}
\end{center}
\end{tcolorbox}

\vspace{5mm}

\begin{multicols}{2}

\subsection*{یک صحنهٔ آشنا}

تصور کنید در یک مهمانی هستید. بحث سیاسی شروع می‌شود. یکی می‌گوید: «این سیاست‌ها کاملاً چپ‌روانه است!» دیگری می‌گوید: «نه، این راست افراطی است!» سومی می‌گوید: «همه‌تان اشتباه می‌کنید، این نئولیبرالیسم است!»

و شما در دل می‌پرسید: این واژه‌ها یعنی چه؟

این صحنه برای من بارها تکرار شده. متوجه شدم که اکثر مردم از این واژه‌ها استفاده می‌کنند بدون اینکه معنای دقیق‌شان را بدانند. «چپ» برای عده‌ای یعنی کمونیسم، برای عده‌ای یعنی لیبرالیسم، برای عده‌ای یعنی هر چیزی که دوست ندارند. همین‌طور «راست».

این سردرگمی فقط در گفتگوهای معمولی نیست. در رسانه‌ها، در تحلیل‌های سیاسی، حتی گاهی در متون آکادمیک، این مفاهیم به شکل مبهم و متناقض به کار می‌روند.

\subsection*{چرا این مهم است؟}

شاید بپرسید: خب چه اهمیتی دارد؟ بگذار هرکس هرچه می‌خواهد بگوید.

جواب این است که سیاست بر زندگی ما اثر می‌گذارد. کیفیت آموزش، بهداشت، امنیت، اقتصاد—همه به تصمیمات سیاسی بستگی دارند. اگر نفهمیم جریان‌های سیاسی چه می‌گویند و چه می‌خواهند، نمی‌توانیم انتخاب آگاهانه کنیم.

بدتر از آن، وقتی نفهمیم «دیگری» چه می‌گوید، او را شیطانی می‌کنیم. تصور می‌کنیم هر که با ما مخالف است، یا احمق است یا شرور. این راه به جایی نمی‌برد.

\subsection*{طیف سیاسی: نقشه‌ای برای سفر}

طیف سیاسی مانند نقشه‌ای است که به ما کمک می‌کند در سرزمین پیچیدهٔ سیاست راه بیابیم. مانند هر نقشه‌ای، این طیف ساده‌سازی است. واقعیت پیچیده‌تر است. اما بدون نقشه، گم می‌شویم.

این کتاب تلاش می‌کند:
\begin{enumerate}
    \item \textbf{تاریخ را روایت کند:} مفاهیم چپ و راست از کجا آمدند و چگونه تحول یافتند.
    \item \textbf{جریان‌ها را معرفی کند:} هر جریان فکری چه می‌گوید و چرا.
    \item \textbf{ارتباطات را نشان دهد:} جریان‌ها چگونه به هم مرتبط‌اند و چه تفاوت‌هایی دارند.
    \item \textbf{ابزار تحلیل بدهد:} چگونه می‌توان رویدادهای جاری را فهمید.
\end{enumerate}

\subsection*{محدودیت‌ها}

این کتاب جامع نیست و نمی‌تواند باشد. تاریخ اندیشهٔ سیاسی دریایی است که عمر انسان برای شنا کردن در همهٔ آن کافی نیست. من انتخاب کرده‌ام که بر جریان‌های اصلی تمرکز کنم و بسیاری از جزئیات را حذف کنم.

همچنین، تمرکز اصلی کتاب بر غرب است. این نه به این دلیل است که غیرغرب مهم نیست، بلکه به این دلیل است که مفاهیم چپ و راست در غرب زاده شده‌اند. فصل ۱۶ به جهان غیرغربی می‌پردازد، اما این نگاهی اجمالی است، نه جامع.

\subsection*{دعوت به سفر}

من از شما دعوت می‌کنم که این کتاب را نه به عنوان حکم نهایی، بلکه به عنوان نقطهٔ شروع بخوانید. با هر فصل کتاب‌هایی معرفی می‌شوند که می‌توانید برای مطالعهٔ بیشتر به آنها مراجعه کنید.

مهم‌تر از همه، از شما می‌خواهم که فکر کنید. با متن موافق باشید یا مخالف، مهم این است که فعالانه درگیر شوید. پرسش‌های تأملی در پایان هر فصل برای همین هدف طراحی شده‌اند.

اکنون سفر آغاز می‌شود. به فرانسهٔ ۱۷۸۹ می‌رویم، جایی که همه چیز شروع شد.

\end{multicols}